%*******************************************************
% 4. Đề cương sơ bộ dự kiến — 4.1 Lý do chọn đề tài
%*******************************************************
\chapter{Đề cương sơ bộ dự kiến}\label{ch:4}

Chương này là phần chính của đề cương. Nó đi theo trình tự của một nghiên cứu thực nghiệm: vì sao câu hỏi đáng hỏi (mục~\ref{sec:lydo}), người khác đã trả lời đến đâu (mục~\ref{sec:tongquan}), đề tài sẽ đo và kiểm định như thế nào (mục~\ref{sec:phuongphap}), đo cái gì ở đâu và trong bao lâu (mục~\ref{sec:phamvi}), và kết quả sẽ được viết thành luận án ra sao (mục~\ref{sec:bocuc}).

%=======================================================
\section{Lý do chọn đề tài}\label{sec:lydo}
%=======================================================

\begin{tomtat}
Hà Nội muốn trở thành trung tâm tài chính chất lượng cao, nhưng trung tâm tài chính quốc tế được đặt ở Thành phố Hồ Chí Minh và Đà Nẵng, và sàn cổ phiếu niêm yết của Hà Nội sắp chuyển hết vào HOSE. Lúc này chưa ai đo được Hà Nội mất hay được gì từ việc đó. Đề tài đo ba thứ: sàn giao dịch liên hệ thế nào với vị thế trung tâm tài chính (câu hỏi~1), sàn HNX kém HOSE bao nhiêu trước hợp nhất (câu hỏi~2), và việc chuyển sàn làm thanh khoản và khả năng đầu tư của các cổ phiếu chuyển sàn thay đổi bao nhiêu (câu hỏi~3). Sản phẩm là hai chỉ số đo được, một tệp đăng ký trước, ba bài báo và một báo cáo định hướng cho các thị trường mà HNX tiếp tục vận hành.
\end{tomtat}

\subsection*{Bài toán số 19 của Hà Nội}

Hà Nội là trung tâm tài chính lớn của cả nước, nhưng hệ thống dịch vụ tài chính chưa đạt tầm khu vực. Thành phố đưa vấn đề này vào Danh mục bài toán lớn giai đoạn 2026--2030 dưới dạng bài toán số 19: \enquote{Phương thức nào để xây dựng Hà Nội thành trung tâm tài chính và dịch vụ chất lượng cao quốc tế?} (Ủy ban nhân dân thành phố Hà Nội, 2026, Phụ lục, bài toán số 19).

Bài toán số 19 được chia thành bốn câu hỏi. Câu hỏi số 1 hỏi những mô hình trung tâm tài chính nào phù hợp với điều kiện phát triển của Hà Nội. Câu hỏi số 2 hỏi về công nghệ tài chính giúp hiện đại hóa hệ thống tài chính đô thị. Câu hỏi số 3 hỏi những nền tảng công nghệ nào giúp tăng cường minh bạch và hiệu quả của thị trường tài chính. Câu hỏi số 4 hỏi cách phát triển các dịch vụ tài chính thế hệ mới dựa trên công nghệ tài chính và kinh tế số. Đi kèm bốn câu hỏi là năm sản phẩm thành phố mong đợi, trong đó có \enquote{Trung tâm tài chính quốc tế tại Hà Nội} và \enquote{Hệ thống dữ liệu tài chính đô thị}.

Đề tài không đề xuất mô hình trung tâm tài chính (câu hỏi số 1 của bài toán) hay nền tảng công nghệ cho thị trường tài chính (câu hỏi số 3). Đề tài cung cấp thước đo nền cho hai câu hỏi này: chất lượng giao dịch cổ phiếu khác nhau thế nào giữa sàn niêm yết của Hà Nội và sàn của Thành phố Hồ Chí Minh, và thay đổi thế nào khi cổ phiếu niêm yết chuyển từ Sở Giao dịch Chứng khoán Hà Nội (HNX) sang Sở Giao dịch Chứng khoán Thành phố Hồ Chí Minh (HOSE).

Lý do chọn vai trò \enquote{thước đo nền} là một luận án tiến sĩ có thể làm tốt việc đo, còn việc chọn mô hình trung tâm tài chính cho một thành phố là quyết định chính sách cần nhiều đầu vào khác. Câu hỏi số 1 khó trả lời nếu chưa biết thị trường hiện có của Hà Nội đang ở mức nào; câu hỏi số 3 khó trả lời nếu chưa có cách đo \enquote{minh bạch và hiệu quả} của thị trường. Đề tài cung cấp đúng hai thứ còn thiếu ấy.

\subsection*{Các khái niệm dùng trong đề cương}

Đề cương dùng các khái niệm sau với nghĩa cố định. \emph{Năng lực cạnh tranh của trung tâm tài chính} là khả năng giữ và thu hút các tổ chức tài chính và cung cấp dịch vụ đáp ứng nhu cầu trong nước và quốc tế (Yao, 2012), đo bằng điểm Chỉ số Trung tâm Tài chính Toàn cầu (GFCI). \emph{Chất lượng thị trường vốn} gồm ba chiều: thanh khoản, hiệu quả thông tin và minh bạch công bố. \emph{Khả năng đầu tư} của một rổ cổ phiếu là lợi suất tương đương chắc chắn sau chi phí giao dịch mà nhà đầu tư đa dạng hóa đạt được với rổ đó, khác với nghĩa giới hạn sở hữu của nhà đầu tư nước ngoài (Bae và cộng sự, 2004).

\emph{Hợp nhất sàn} là việc chuyển toàn bộ cổ phiếu niêm yết từ HNX sang HOSE. \emph{Thị trường vốn Hà Nội} là các thị trường do HNX tiếp tục vận hành sau hợp nhất: trái phiếu, chứng khoán phái sinh, và thị trường giao dịch cổ phiếu của công ty đại chúng chưa niêm yết (UPCoM) khi thị trường này còn do HNX tổ chức.

\begin{delam}
\emph{Dễ nghĩ là} \enquote{khả năng đầu tư} ở đây cũng là chữ \emph{investability} quen thuộc trong tài chính quốc tế, tức nhà đầu tư nước ngoài được phép mua bao nhiêu phần trăm cổ phần. Không phải vậy. Trong đề cương, khả năng đầu tư đo \emph{nhà đầu tư thu được bao nhiêu} khi đa dạng hóa trên một rổ cổ phiếu, sau khi trừ chi phí giao dịch và tính cả rủi ro. Hai rổ có cùng mức room ngoại vẫn có thể có khả năng đầu tư rất khác nhau nếu một rổ khó mua bán hơn.
\end{delam}

\subsection*{Vị thế hiện tại của Hà Nội}

Theo Nghị quyết số 222/2025/QH15, Trung tâm tài chính quốc tế tại Việt Nam đặt tại Thành phố Hồ Chí Minh và thành phố Đà Nẵng (Quốc hội, 2025). Ở kỳ GFCI 40 (tháng 9/2026), Thành phố Hồ Chí Minh xếp thứ 67 trên 117 trung tâm, Đà Nẵng lần đầu vào chỉ số ở vị trí 71, còn Hà Nội không có trong bảng xếp hạng (Wardle \& Mainelli, 2026).

Hai sự kiện này đặt Hà Nội vào một vị trí khó: thủ đô có quy mô tài chính lớn nhưng không có trung tâm tài chính quốc tế theo luật, không có mặt trong chỉ số quốc tế, và sắp không còn sàn cổ phiếu niêm yết. Câu hỏi thực tế là Hà Nội còn những lợi thế nào trên thị trường vốn và nên dồn sức vào đâu.

\subsection*{Lộ trình hợp nhất}

Lộ trình hợp nhất đã bị lùi hai lần: hạn chuyển cổ phiếu niêm yết từ HNX sang HOSE là năm 2023 theo Thông tư số 57/2021/TT-BTC, hết năm 2025 theo Thông tư số 69/2023/TT-BTC (Bộ Tài chính, 2023) và hết ngày 31/12/2026 theo Thông tư số 139/2025/TT-BTC (Bộ Tài chính, 2025). Theo thông báo ngày 30/9/2026 của Sở Giao dịch Chứng khoán Việt Nam (VNX), cổ phiếu niêm yết tại HNX giao dịch phiên cuối ngày 23/12/2026, tạm ngừng hai ngày 24 và 25/12/2026, và giao dịch tại HOSE từ ngày 28/12/2026 (Sở Giao dịch Chứng khoán Việt Nam, 2026).

Hạn chuyển cổ phiếu UPCoM sang HOSE cũng là hết ngày 31/12/2026, do Thông tư số 69/2023/TT-BTC đặt và không bị Thông tư số 139/2025/TT-BTC sửa (Bộ Tài chính, 2023, 2025). Tuy vậy, thông báo của VNX chỉ áp dụng cho cổ phiếu niêm yết; báo chí cho biết cổ phiếu UPCoM không thuộc diện chuyển trong đợt này (Báo Đầu tư Chứng khoán, 2026), và đến ngày 17/10/2026 chưa có văn bản công bố lịch chuyển cổ phiếu UPCoM; đề tài sẽ cập nhật lịch này khi VNX công bố (mục~\ref{sec:tincay}).

Việc lộ trình bị lùi hai lần có hai hệ quả cho thiết kế nghiên cứu. Thứ nhất, thị trường đã biết trước khả năng hợp nhất từ năm 2021, nên một phần phản ứng có thể đã xảy ra trước ngày chuyển sàn; đề tài vì vậy tách riêng giai đoạn chuyển tiếp và kiểm tra phản ứng trước sự kiện (mục~\ref{sec:kiemdinh}). Thứ hai, ngày chuyển sàn vẫn có thể thay đổi; đề tài chuẩn bị sẵn phương án cho trường hợp chuyển theo nhiều đợt hoặc bị lùi tiếp (mục~\ref{sec:tincay}). Các mốc pháp lý được gom lại ở Phụ lục~\ref{pl:mocthoigian}.

\subsection*{Lý thuyết nói gì về nhiều trung tâm trong một nước}

Lý thuyết về trung tâm tài chính giải thích vì sao hoạt động tài chính dồn về một vài thành phố: quy mô giao dịch và thông tin tạo lợi thế tích tụ (Gehrig, 2000; Kindleberger, 1974). Khi một quốc gia có nhiều trung tâm, các trung tâm phân vai: ở Trung Quốc, Thượng Hải là trung tâm thương mại, Bắc Kinh là trung tâm chính trị và quản lý, Hồng Kông là trung tâm hải ngoại (Lai, 2012), và Bắc Kinh là trung tâm tài chính lớn dù không có sở giao dịch chứng khoán tầm quốc gia trong phần lớn quá trình phát triển (Naeem, 2023). Với Hà Nội, câu hỏi sau hợp nhất là nên chuyên môn hóa vào những thị trường nào, tương tự cách Yao (2012) phân tích lợi thế của Luxembourg đối với ngành quỹ đầu tư.

Trường hợp Bắc Kinh cho thấy mất sàn cổ phiếu chưa chắc là mất vị thế tài chính: một thủ đô có thể giữ vai trò lớn nhờ cơ quan quản lý, trụ sở ngân hàng và thông tin chính sách. Trường hợp Luxembourg cho thấy một trung tâm nhỏ vẫn cạnh tranh được nếu chọn đúng mảng chuyên môn. Đề tài không giả định Hà Nội sẽ đi theo hướng nào; đề tài đo những gì cần đo để lựa chọn có căn cứ.

\subsection*{Vì sao đây là một cơ hội đo hiếm có}

Đợt hợp nhất ở Việt Nam đổi sàn của toàn bộ cổ phiếu niêm yết tại HNX, gồm 299 mã tính đến phiên 30/9/2026 trong cùng một ngày (Tuổi Trẻ Online, 2026b). Sự kiện này cho phép so sánh cùng các cổ phiếu trước và sau khi đổi sàn với một nhóm đối chứng, về thanh khoản lẫn khả năng đầu tư.

Trong nghiên cứu thông thường về sàn giao dịch, doanh nghiệp tự chọn sàn. Doanh nghiệp chuyển sang sàn lớn thường là doanh nghiệp đang tăng trưởng, nên thanh khoản tăng sau khi chuyển có thể do doanh nghiệp tốt lên chứ không do sàn. Ở đợt hợp nhất năm 2026, doanh nghiệp không chọn: mọi mã niêm yết tại HNX đều chuyển, cùng một ngày, theo quyết định hành chính. Đây là điều kiện gần với thí nghiệm nhất mà nghiên cứu về sàn giao dịch có thể có.

\subsection*{Ba vấn đề tài liệu còn để ngỏ}

Tài liệu hiện có để ngỏ ba vấn đề. Thứ nhất, 17 yếu tố mà Naeem (2023) dùng để giải thích thứ hạng GFCI không có chỉ báo nào về thị trường cổ phiếu, còn các nghiên cứu khác xét nhiều nhóm yếu tố cùng lúc nên chưa tách riêng vai trò của thị trường cổ phiếu (Eichengreen \& Shah, 2020; Moosa và cộng sự, 2016). Thứ hai, bằng chứng về hợp nhất sàn chủ yếu đến từ thị trường phát triển, trái chiều theo quy mô doanh nghiệp và chỉ đo thanh khoản (Hisada, 2018; Nielsson, 2009). Thứ ba, các công trình về Việt Nam tập trung vào HOSE (Goutte và cộng sự, 2025; Hoang \& Luu, 2024; Vo \& Phan, 2017) và chưa có nghiên cứu nào về đợt hợp nhất năm 2026.

\subsection*{Câu hỏi nghiên cứu}

Từ những vấn đề trên, đề tài đặt ra ba câu hỏi, mỗi câu gắn với câu hỏi tương ứng của bài toán số 19 (ghi trong ngoặc):

\begin{description}
\item[Câu hỏi 1:] Mức phát triển thị trường cổ phiếu liên hệ thế nào với điểm GFCI của trung tâm tài chính, sau khi kiểm soát chất lượng thể chế và thu nhập? (bối cảnh cho câu hỏi số 1 của bài toán số 19)
\item[Câu hỏi 2:] Trước khi hợp nhất, chất lượng thị trường vốn của sàn niêm yết HNX thấp hơn của HOSE bao nhiêu? (câu hỏi số 3)
\item[Câu hỏi 3:] Việc chuyển sàn ngày 28/12/2026 làm thanh khoản và khả năng đầu tư của các cổ phiếu chuyển sàn thay đổi bao nhiêu? (câu hỏi số 1 và số 3)
\end{description}

Ba câu hỏi đi từ rộng đến hẹp. Câu hỏi 1 nhìn ra thế giới: ở các trung tâm tài chính khác, thị trường cổ phiếu phát triển có đi cùng vị thế cao hơn không. Câu hỏi 2 nhìn vào Việt Nam trước sự kiện: hai sàn khác nhau đến đâu. Câu hỏi 3 nhìn vào chính sự kiện: khi cùng một cổ phiếu đổi sàn, điều gì thay đổi. Câu trả lời của câu hỏi sau có ý nghĩa hơn khi đã có câu trả lời của câu hỏi trước.

\subsection*{Mục tiêu nghiên cứu}

Mục tiêu tổng quát là đo khoảng cách chất lượng thị trường giữa hai sàn trước hợp nhất và lượng hóa tác động của việc chuyển sàn năm 2026 lên thanh khoản và khả năng đầu tư của các cổ phiếu chuyển sàn, làm căn cứ cho định hướng chuyên môn hóa của thị trường vốn Hà Nội.

Bốn mục tiêu cụ thể là:
\begin{enumerate}[label=(\roman*)]
\item ước lượng liên hệ giữa mức phát triển thị trường cổ phiếu và điểm GFCI trên dữ liệu bảng quốc tế;
\item xây dựng Chỉ số chất lượng thị trường vốn theo tháng cho sàn niêm yết HNX, UPCoM và HOSE;
\item ước lượng tác động của việc chuyển sàn lên thanh khoản và khả năng đầu tư, kể cả khác biệt của tác động theo mức chịu cú sốc biên độ, bước giá, chỉ số (mục~\ref{sec:kiemdinh}) và theo trụ sở doanh nghiệp;
\item tổng hợp định tính các kết quả thành định hướng chuyên môn hóa và hạ tầng dữ liệu cho thị trường vốn Hà Nội.
\end{enumerate}

Trong bốn giả thuyết ở mục~\ref{sec:giathuyet}, \HH{3} là kiểm định trọng tâm; \HH{1} cung cấp bối cảnh quốc tế, \HH{2} cung cấp thước đo nền trước sự kiện, và \HH{4} quy đổi tác động sang khả năng đầu tư.

\begin{delam}
\emph{Dễ nghĩ là} bốn giả thuyết có vai trò ngang nhau. Thực ra chỉ \HH{3} là kiểm định trung tâm, vì đó là chỗ đề tài có thiết kế nhận dạng mạnh nhất (một sự kiện ngoại sinh, có nhóm đối chứng). \HH{1} chỉ cho kết luận về liên hệ có điều kiện, không phải quan hệ nhân quả. \HH{2} so sánh hai sàn có nhiều điểm khác nhau ngoài bản thân sàn. \HH{4} chỉ được ước lượng vì có quá ít cửa sổ để kiểm định có lực. Khi trình bày, nên giữ đúng thứ bậc này.
\end{delam}

\subsection*{Đóng góp dự kiến}

So với Yao (2012), Eichengreen và Shah (2020) và Naeem (2023), đề tài đưa riêng mức phát triển thị trường cổ phiếu vào mô hình điểm GFCI. So với Nielsson (2009) và Hisada (2018), đề tài nghiên cứu hợp nhất hai sàn trong một quốc gia mới nổi, dùng suy luận phù hợp với một sự kiện duy nhất và đo thêm tác động lên khả năng đầu tư. So với các công trình về Việt Nam (Goutte và cộng sự, 2025; Nguyễn \& Nguyễn, 2026a; Vo \& Phan, 2017), đề tài xây dựng chỉ số so sánh HNX với HOSE và cung cấp bằng chứng đầu tiên, trong phạm vi tìm kiếm của thí sinh, về đợt hợp nhất năm 2026.

Về sản phẩm, Chỉ số chất lượng thị trường vốn có thể là một cấu phần của sản phẩm \enquote{Hệ thống dữ liệu tài chính đô thị} mà bài toán số 19 nêu, còn báo cáo định hướng chuyên môn hóa là một đầu vào phân tích cho sản phẩm \enquote{Trung tâm tài chính quốc tế tại Hà Nội}. Đề tài không thiết kế sản phẩm \enquote{Trung tâm tài chính quốc tế tại Hà Nội}, và Nghị quyết số 222/2025/QH15 hiện chỉ đặt trung tâm tài chính quốc tế tại Thành phố Hồ Chí Minh và Đà Nẵng. Hai chỉ số (chất lượng thị trường vốn và khả năng đầu tư) cùng tệp đăng ký trước là sản phẩm có được dù kết quả kiểm định ra sao.

Điểm cuối cùng quan trọng với người cấp kinh phí. Một đề tài chỉ có giá trị khi giả thuyết được ủng hộ sẽ khuyến khích nhà nghiên cứu tìm cách cho giả thuyết được ủng hộ. Ở đây, ba sản phẩm (hai chỉ số và tệp đăng ký) có giá trị độc lập với chiều của kết quả, nên thành phố nhận được sản phẩm dùng được dù việc chuyển sàn làm thanh khoản tăng, giảm hay không đổi.

\subsection*{Sáu chỉ tiêu kiểm chứng}

Kết quả của đề tài được kiểm chứng bằng sáu chỉ tiêu; mỗi chỉ tiêu có mốc thời gian và bằng chứng công khai:

\begin{enumerate}[label=(\arabic*)]
\item Tệp đăng ký trước kế hoạch phân tích trên OSF Registries, khóa có dấu thời gian trước ngày 28/12/2026.
\item Chỉ số chất lượng thị trường vốn theo tháng cho ba sàn từ tháng 6/2009, công bố cùng mã nguồn trên kho lưu trữ mở.
\item Ước lượng tác động của việc chuyển sàn lên thanh khoản với 24 tháng dữ liệu sau sự kiện, kèm hiệu ứng nhỏ nhất phát hiện được và khoảng tin cậy, theo đúng tệp đăng ký.
\item Chỉ số khả năng đầu tư của rổ cổ phiếu chuyển sàn, cho 20 cửa sổ trước và 8 cửa sổ sau sự kiện.
\item Ba bài báo trên tạp chí có phản biện, với mục tiêu ít nhất một bài, do thí sinh là tác giả chính, thuộc cơ sở dữ liệu Web of Science hạng Q1 hoặc Q2, và một đến hai báo cáo hội thảo đăng trong kỷ yếu được Hội đồng Giáo sư nhà nước tính điểm (hội thảo trong nước) hoặc thuộc cơ sở dữ liệu Scopus hoặc Web of Science (hội thảo ở nước ngoài).
\item Báo cáo định hướng chuyên môn hóa cho các thị trường do HNX tiếp tục vận hành, gửi Sở Khoa học và Công nghệ Hà Nội trong năm thứ tư, kèm văn bản xác nhận hoặc cam kết sử dụng kết quả của ít nhất một đơn vị, ưu tiên HNX.
\end{enumerate}

Mỗi chỉ tiêu có thể được người ngoài kiểm tra mà không cần tin lời thí sinh: tệp đăng ký có dấu thời gian, chỉ số và mã nguồn công khai, bài báo có DOI, báo cáo có văn bản tiếp nhận. Chỉ tiêu (1) có hạn sớm nhất và quan trọng nhất, vì nếu tệp đăng ký được khóa sau ngày 28/12/2026 thì không còn chứng minh được rằng kế hoạch phân tích có trước dữ liệu.
